\section{Semipositive metrics in a four-point blowup family}
\label{sec:four-point-metric-nonextension}\index{Semipositive metric!failure of extension}\index{Blowup!four-point family}

A semipositive metric on a special fiber need not be the restriction of a semipositive metric on the total space. The obstruction below comes from compact curves on the nearby fibers: degree zero forces their restricted curvature to vanish, and this condition persists on a component of the central curve. The prescribed metric has a smaller logarithmic coefficient than this condition permits. Varying these curves will also identify every central metric that does extend.

Let $\Delta$ be a disk with coordinate $t$. Blow up $\mathbb P^2\times\Delta$ along the four disjoint sections
\[
 \begin{aligned}
 p_1(t)&=[1:0:0],&\qquad p_2(t)&=[0:1:0],\\
 p_3(t)&=[1:1:t],&\qquad p_4(t)&=[0:0:1].
 \end{aligned}
\]
and write $\pi:\mathcal X\to\Delta$ for the resulting family and $X_t=\pi^{-1}(t)$. Let $H$ be the pulled-back hyperplane line bundle, let $E_i$ be the exceptional divisors, and put
\[
 L=2H-E_1-E_2-E_3-E_4,\qquad L_t=L|_{X_t}.
\]
We use additive notation for tensor products of line bundles and also for their divisor classes, and write $K_{X_t}$ for the canonical line bundle. A singular Hermitian metric has squared norm $|e|_h^2=\mathrm{e}^{-\varphi}$ in a local holomorphic frame; its curvature is $\ddc\varphi$. Thus the divisor metric has weight $\log|w|^2$ and curvature $[w=0]$.

\begin{lemma}\label{lem:four-point-family-adjoint-vanishing}
The family is smooth and projective with rational surface fibers, and $L$ is effective on $\mathcal X$. For every $t\ne0$, $L_t$ is generated by global sections, has square zero and is nef. For every integer $m\geq1$,
\[
 H^0(X_t,K_{X_t}+mL_t)=0\qquad(t\ne0).
\]
\end{lemma}
\begin{proof}
Near a section, coordinates centered on it identify the blowup with $\operatorname{Bl}_0(\mathbb C^2)\times\Delta$. The four centers are disjoint, so these charts prove smoothness and the submersion property. The blowup of the projective morphism is projective and proper, and every fiber is the indicated rational surface.

The quadratic forms
\[
 F_1=x(z-ty),\qquad F_2=y(z-tx)
\]
vanish along all four sections and lift to sections of $L$. Each is nonzero, proving effectivity. For $t\ne0$ their common zeros in $\mathbb P^2$ are exactly the four points. Centered affine coordinates at $p_1,p_2,p_3,p_4$ give the linear initial terms
\[
 (v-tu,-tu),\quad(-tu,v-tu),\quad(v-tu,v),\quad(u,v).
\]
Each pair is independent. After subtracting the exceptional divisors, no common zero remains on an exceptional curve. The two transformed sections therefore generate $L_t$, whose degree on every irreducible curve is nonnegative.

On a fiber, $H^2=1$, $H\cdot E_i=0$, and $E_i\cdot E_j=-\delta_{ij}$. Hence $L_t^2=0$ and $K_{X_t}=-3H+\sum_iE_i$ gives $K_{X_t}\cdot L_t=-2$. A nonzero adjoint section would have an effective zero divisor with negative intersection against the nef line $L_t$, a contradiction. This also excludes a nowhere-vanishing section, whose zero divisor has intersection zero.
\end{proof}

On $X_0$, write $\ell$ for the strict transform of $z=0$, and put
\[
 \ell=H-E_1-E_2-E_3,\qquad
 A=2L_0-\ell=3H-E_1-E_2-E_3-2E_4.
\]

\begin{lemma}\label{lem:four-point-central-semipositive-metric}
The line bundle $A$ is generated by global sections. There is a semipositive metric $h_0$ on $L_0$ whose weight along $\ell$ is
\[
 \varphi_0=\tfrac12\log|w|^2+\text{a smooth function},
\]
where $w$ is a local equation of $\ell$. It is smooth off $\ell$ and has analytic singularities.
\end{lemma}
\begin{proof}
Use the cubic sections
\[
 x^2z,\qquad xyz,\qquad y^2z,\qquad xy(x-y).
\]
Their common zero set in $\mathbb P^2$ is the four central points. At $p_1,p_2,p_3$, their linear initial terms include respectively $(z,y)$, $(z,x)$ and $(z,x-y)$. At $p_4$, the quadratic initial terms include $x^2,xy,y^2$. These have no common zero on the respective exceptional projective lines. Subtracting the specified simple and double exceptional factors therefore gives generating sections of $A$.

Their Fubini--Study metric $h_A$ is smooth and semipositive. Let $h_\ell$ be the divisor metric on $\mathcal O(\ell)$. Since $2L_0=A+\ell$, the square root of $h_Ah_\ell$ defines $h_0$, with
\[
 \Theta_{h_0}=\tfrac12\Theta_{h_A}+\tfrac12[\ell]\geq0.
\]
The displayed weight and its asserted singularities follow directly.
\end{proof}

\begin{lemma}\label{lem:four-point-conic-degeneration}\index{Curvature!degree-zero curve}
The strict transforms of the conics
\[
 F_1-2F_2=z(x-2y)+txy=0
\]
form a surface $\Gamma\subset\mathcal X$. For $t\ne0$, its fiber $C_t$ is a smooth compact rational curve with $L_t\cdot C_t=0$. The central fiber is $\ell+q$, where $q$ is the strict transform of $x=2y$. The curves $\ell,q$ meet transversely at one point $p$, and $L_0\cdot q=1$. Near every point of $q\setminus\{p\}$, $\Gamma$ is smooth and $\pi|_\Gamma$ is a submersion.
\end{lemma}
\begin{proof}
The symmetric matrix of the quadratic form has determinant $-t/2$, so its conic is smooth for $t\ne0$. It passes simply through all four centers; its strict transform is still a smooth rational curve of class $L_t$. Thus its degree against $L_t$ is $L_t^2=0$.

At $t=0$ the conic is the union of $z=0$ and $x=2y$. The latter avoids $p_1,p_2,p_3$ and passes through $p_4$. Both components have multiplicity one at each center they meet, so the central transformed divisor has no extra exceptional component. Their intersection $[2:1:0]$ is away from every center and transverse. Since $q=H-E_4$, one has $L_0\cdot q=1$.

Away from the centers, the central conic equation has nonzero spatial derivative at every point of $q\setminus\{p\}$, proving the submersion claim. Near $p_4$, use $z=1$, $x=r$, $y=rs$ in a blowup chart. After removing the exceptional factor, the equation is
\[
 1-2s+trs=0.
\]
At the central point on $q\cap E_4$, $r=0$ and $s=1/2$, and its derivative in $s$ is $-2$. Thus $(r,t)$ are local coordinates on $\Gamma$ there. These checks cover $q\setminus\{p\}$.
\end{proof}

\begin{theorem}\label{thm:four-point-semipositive-metric-nonextension}\index{Singularity type!metric extension}
There is no semipositive metric on $L$ over any open neighborhood of $X_0$ whose restricted weights $\psi$ on $X_0$ are well defined and satisfy
\[
 \psi\geq\varphi_0-C
\]
in compatible local frames, for a constant $C$. In particular, $h_0$ has neither an exact semipositive extension nor one of the same singularity type.
\end{theorem}
\begin{proof}
Suppose such a metric exists. Properness lets us shrink $\Delta$ so that the whole restricted family lies in its domain: a sequence outside the neighborhood over $t_j\to0$ would have a limit outside it on $X_0$. The comparison makes its weight finite at every point of $q\setminus\{p\}$, because $\varphi_0$ is smooth there. Hence restriction to $\Gamma$ and to $q$ is well defined.

Choose a product chart $(\zeta,t)$ on $\Gamma$ near a point of $q\setminus\{p\}$ and a holomorphic frame of $L$. Its restricted weight $u(\zeta,t)$ is plurisubharmonic and locally integrable. Fubini implies that, for almost every small $t\ne0$, the metric restricts genuinely to $C_t$. Its positive curvature measure on this compact curve has total mass $L_t\cdot C_t=0$, so it vanishes. Thus $u(\cdot,t)$ is harmonic for almost every $t$.

We check that this also forces the central slice to be harmonic. Distributional differentiation in $\zeta$ and Fubini give $\partial_\zeta\partial_{\bar\zeta}u=0$ on the product chart. For small $\varepsilon>0$, the disk averages
\[
 u_\varepsilon(\zeta)=\frac1{\pi\varepsilon^2}
       \int_{|t|<\varepsilon}u(\zeta,t)\,\mathrm{d}A(t)
\]
use planar area $\mathrm{d}A$. They are finite, since subharmonicity in $t$ bounds them below by the finite central weight, and local upper bounds for $u$ bound them above. They are subharmonic in $\zeta$ and have zero distributional Laplacian, so their canonical representatives are harmonic. Subharmonicity in $t$ also makes these averages decrease to $u(\zeta,0)$ as $\varepsilon\downarrow0$. A decreasing sequence of harmonic functions has a harmonic limit or tends identically to $-\infty$: Harnack's inequality applied after subtracting from the first function bounds the differences on compact subdisks whenever their values are bounded at one point, and harmonic compactness gives the limit. Here the central weight is finite, so the limit is harmonic. This argument applies at every point of $q\setminus\{p\}$.

The restricted metric on $L_0|_q$ therefore has curvature supported at $p$. Positivity and $L_0\cdot q=1$ make this measure $[p]$. In a coordinate $w$ on $q$ centered at $p$, its weight is $\log|w|^2$ plus a harmonic function, since subtracting this logarithm has zero distributional Laplacian. Transversality gives $\varphi_0|_q=\tfrac12\log|w|^2+O(1)$. Their difference tends to $-\infty$, contradicting the assumed lower bound. This proves the theorem.
\end{proof}

The obstruction concerns the prescribed singularity type. Since $L$ is effective on the total space, a divisor metric supplies a different semipositive metric there.

\subsection{The image of metric restriction}
\label{subsec:four-point-metric-restriction-image}\index{Semipositive metric!restriction image}\index{Pencil!metric extension}

The degree-zero conics also determine which central metrics extend. A restricted metric is called genuine if its local weights are not identically $-\infty$.

\begin{lemma}\label{lem:psh-harmonic-slice-persistence}
Let $u$ be a genuine plurisubharmonic function on a product of disks $D_\zeta\times D_t$, with $0\in D_t$. Suppose $u(\cdot,t)$ is harmonic for almost every $t$. Then $u(\cdot,0)$ is harmonic or identically $-\infty$.
\end{lemma}
\begin{proof}
Local integrability and Fubini give $\partial_\zeta\partial_{\bar\zeta}u=0$ in distributions. On smaller product disks, form the averages
\[
 u_\varepsilon(\zeta)=\frac1{\pi\varepsilon^2}
       \int_{|t|<\varepsilon}u(\zeta,t)\,\mathrm{d}A(t),
\]
where $\mathrm{d}A$ is planar area. Local upper bounds and reverse Fatou give upper semicontinuity, and Fubini gives the submean inequality in $\zeta$. The averages are locally integrable, since $u$ is integrable on compact product disks. Thus each $u_\varepsilon$ is subharmonic with zero distributional Laplacian, and is harmonic in its canonical representative.

Subharmonicity in $t$ makes these disk averages decrease to $u(\zeta,0)$ as $\varepsilon\downarrow0$. The decreasing limit of harmonic functions is harmonic or identically $-\infty$: subtract the functions from the first one and apply Harnack's inequality and harmonic compactness whenever the values are bounded at one point. This proves the alternative on every smaller disk and hence on $D_\zeta$.
\end{proof}

Let $f:X_0\to\mathbb P^1$ be the resolved pencil of lines through $p_4$. Its affine coordinate is $a=x/y$, and put $P=f^*\mathcal O_{\mathbb P^1}(1)$.

\begin{lemma}\label{lem:four-point-pencil-reference}
The pencil is a holomorphic map with connected fibers, and $P=H-E_4$ and $L_0=\mathcal O(\ell)+P$. Let $h_\ell$ be the divisor metric normalized by $|s_\ell|_{h_\ell}=1$ away from $\ell$, where $s_\ell$ is its canonical section, and let $h_{\mathrm{FS}}$ be the Fubini--Study metric on $\mathcal O_{\mathbb P^1}(1)$, with curvature $\omega_{\mathrm{FS}}$. The metric
\[
 h_*=h_\ell\otimes f^*h_{\mathrm{FS}}
\]
is a genuine restriction of a semipositive metric on $L$ over the whole family. Its curvature is $[\ell]+f^*\omega_{\mathrm{FS}}$.
\end{lemma}
\begin{proof}
First blow up $p_4$ in $\mathbb P^2$. The map $[x:y]$ becomes holomorphic: in the chart $x=r$, $y=rs$, it is $[1:s]$, and the other chart covers the missing direction. Its fibers are the strict lines through $p_4$ and its pencil line is $H-E_4$. Blow up the other three central points. Their directions are $\infty,0,1$, respectively. Each affected fiber becomes the union of its strict transform and one exceptional line meeting it; every other fiber remains a projective line. This proves connectedness and the asserted pullback line bundle. The equality $L_0=\mathcal O(\ell)+P$ follows from the divisor classes.

In this tensor identification, the restrictions of $F_1,F_2$ are $s_\ell f^*x,s_\ell f^*y$. Indeed, at $t=0$ their original forms are $xz,yz$; the pullback of $z=0$ is $\ell+E_1+E_2+E_3$, and removing the four exceptional factors leaves exactly the indicated sections. In a local frame of $L$, the weight $\log(|F_1|^2+|F_2|^2)$ defines a semipositive metric: frame changes multiply both coefficients by the same nonvanishing holomorphic function. The sections have no common zero on any nonzero fiber, so the metric is genuine. Its central weight is the sum of the divisor weight and the pulled-back Fubini--Study weight, proving the claim and the curvature formula.
\end{proof}

\begin{lemma}\label{lem:four-point-variable-conics}
For $a\in\mathbb C\setminus\{0,1\}$, the transformed conics
\[
 F_1-aF_2=z(x-ay)+t(a-1)xy=0
\]
have smooth rational fibers of $L_t$-degree zero when $t\ne0$. Their central fiber is $\ell+q_a$, where $q_a=f^{-1}(a)$ is a smooth projective line of $L_0$-degree one. The two components meet transversely at $p_a=[a:1:0]$, away from the blowup centers. Along $q_a\setminus\{p_a\}$, the conic surface is smooth and its projection to the disk is a submersion.
\end{lemma}
\begin{proof}
The symmetric matrix of the displayed quadratic form has determinant $-a(a-1)t/4$. Its conic is therefore smooth for $t\ne0$ and passes simply through the four centers. Its transformed class is $L_t$, so its degree is $L_t^2=0$. At $t=0$ it splits as the stated two lines; $a\notin\{0,1\}$ and its finiteness exclude the other three central centers from $q_a$. Thus no extra exceptional component occurs, and $q_a=H-E_4$ has degree one against $L_0$. Their intersection is transverse and away from all centers.

Off the centers, the central equation has nonzero spatial derivative along $q_a\setminus\{p_a\}$. At its point on $E_4$, the chart $z=1$, $x=r$, $y=rs$ gives, after removing the exceptional factor,
\[
 1-as+t(a-1)rs=0.
\]
At $r=t=0$, $s=1/a$, its derivative in $s$ is $-a\ne0$. These charts establish the final assertion.
\end{proof}

\begin{theorem}\label{thm:four-point-metric-restriction-image}
The semipositive metrics on $L_0$ that arise as genuine restrictions of semipositive metrics on $L$ over a neighborhood of $X_0$ are exactly
\[
 h_\ell\otimes f^*h_B,
\]
where $h_B$ ranges over all semipositive singular Hermitian metrics on $\mathcal O_{\mathbb P^1}(1)$. Every such central metric extends over the whole family.
\end{theorem}
\begin{proof}
For necessity, shrink the disk by properness so that the extending metric is defined on the whole restricted family, and call its genuine central restriction $h$. On the smooth locus of $f$, local integrability and Fubini show that $h|_{q_a}$ is genuine for almost every $a\in\mathbb C\setminus\{0,1\}$. Fix such an $a$. On each local product chart of the conic surface along $q_a\setminus\{p_a\}$, the extending weight is genuine: its central slice, a restriction of the genuine curve metric, is not identically $-\infty$. Fubini gives genuine metric restrictions on almost every nearby compact conic. Their positive curvature measures have degree zero and hence vanish. \Cref{lem:psh-harmonic-slice-persistence} therefore makes the central weight harmonic on each chart. Its identically $-\infty$ alternative is excluded by the genuine central curve restriction.

The curvature of $h|_{q_a}$ is supported at $p_a$. It is positive of degree one, so it equals $[p_a]$. This is also the curvature of $h_*|_{q_a}$. Their weight difference is a globally defined harmonic function on the projective line: locally at $p_a$, subtracting $\log|w|^2$ from either weight leaves a harmonic function, so their difference extends across that point. It is consequently constant. Thus the weight difference $u$ between $h$ and $h_*$ is constant along almost every fiber, away from $\ell$.

Local holomorphic sections of $f$ avoiding $\ell$ and the exceptional curves exist over every base point. On the finite chart use $[a:1:1]$, and near infinity use $[1:b:1]$ with $b=1/a$. In a local base frame $e_B$, the nonvanishing holomorphic section $s_\ell\otimes f^*e_B$ is a frame of $L_0$ along these sections. Composing the weight of $h$ in this frame with a section defines a subharmonic base weight. It is genuine because the weights on the good fibers are harmonic off their node. Two choices agree for almost every base point because $u$ is constant on those fibers; canonical subharmonic representatives then agree everywhere. Base-frame changes give the metric transition law. These weights therefore define a semipositive metric $h_B$ on $\mathcal O_{\mathbb P^1}(1)$. Fubini on local product charts, with the analytic exceptional fibers and $\ell$ having zero ambient volume, shows that $h$ and $h_\ell\otimes f^*h_B$ agree almost everywhere on $X_0$. Their canonical local psh weights therefore agree everywhere. This proves necessity, including the three exceptional fibers.

For sufficiency, represent $h_B$ by its homogeneous weight $V$ on $\mathbb C^2\setminus\{0\}$. In the chart $Y\ne0$, if $\phi_B$ is the weight in the frame given by the section $Y$, set
\[
 V(X,Y)=\phi_B(X/Y)+\log|Y|^2.
\]
The other chart gives the same function by the metric transition law. It is plurisubharmonic and satisfies $V(\lambda X,\lambda Y)=V(X,Y)+\log|\lambda|^2$ for $\lambda\ne0$. The weight difference between $h_B$ and $h_{\mathrm{FS}}$ is bounded above on the compact base, so
\[
 V(X,Y)\leq\log(|X|^2+|Y|^2)+C.
\]
Thus $V$ extends plurisubharmonically to the origin with value $-\infty$, by \cref{lem:qpsh-removable-analytic}. In each local frame of $L$, the weight $V(F_1,F_2)$ is psh and transforms as a metric weight by homogeneity. It is genuine: the two independent generating sections on each nonzero fiber give a nonconstant map onto $\mathbb P^1$, where $h_B$ is genuine. On the central fiber, the factorization in \cref{lem:four-point-pencil-reference} and homogeneity give exactly the divisor weight plus the pulled-back base weight. This is the required extension over $\mathcal X$.
\end{proof}

\begin{corollary}\label{cor:four-point-least-singular-extendable}\index{Singularity type!least singular extendable metric}
The metric $h_*$ has least singular weights among all genuine extendable restrictions: every such restricted weight $\psi$ satisfies $\psi\leq\psi_*+C$ in compatible local frames, where $\psi_*$ denotes the weight of $h_*$. It is itself extendable. This assertion concerns extendable restrictions, rather than all semipositive metrics on $L_0$.
\end{corollary}
\begin{proof}
By \cref{thm:four-point-metric-restriction-image}, write the metric as $h_\ell\otimes f^*h_B$. Its weight difference from $h_*$ is the pullback of the global quasi-psh difference between $h_B$ and $h_{\mathrm{FS}}$. This function is bounded above on $\mathbb P^1$, proving the inequality. Extendability of $h_*$ is \cref{lem:four-point-pencil-reference}.
\end{proof}
